% ch-inputformat.tex -- the input file: TOML, type conversion, map of the blocks
\chapter{The input file}
\label{ch:inputformat}

A Yuzen run is described by exactly one text file, which you pass to the
executable on the command line. That file is a \textbf{TOML} document. This
chapter explains how Yuzen uses TOML: how the block structure of a case maps
onto TOML tables, how a written value becomes a number inside the solver, how
defaults and required keys are expressed, and which mistakes the reader will
\emph{not} tell you about. The keys themselves are documented block by block in
\cref{ch:ref-domain,ch:ref-physics,ch:ref-bcout,ch:ref-numerics}; the map in
\cref{sec:inputformat:map} tells you which of those chapters to open.

Read this chapter once before you write your first input file by hand. Most of
the ways a Yuzen run goes wrong quietly --- rather than aborting with a message
--- are properties of the reader described here, not of the physics.

\section{TOML, chosen by the file name}
\label{sec:inputformat:formats}

The reader layer is \code{yuzen::config::Block}, and it picks the parser from
the file-name extension alone. A name ending in \file{.toml} is parsed with
toml\texttt{++}; that is the format this manual specifies, and the one to write
new work in. Nothing inside the file influences the choice, so renaming a TOML
file to something else does not convert it --- it hands the file to the \file{.inp}
parser of \cref{sec:inputformat:legacy}, which then fails on TOML syntax.

TOML itself is a published standard and is not re-specified here. Its grammar
--- comments, strings, numbers, booleans, arrays, tables --- is at
\url{https://toml.io}, and any TOML parser, editor mode or scripting library
will read a Yuzen input file correctly. What is specific to Yuzen, and what the
rest of this chapter documents, is which TOML constructs the solver's block
tree is built from, and what happens to a value on its way from the file into a
variable.

\section{Blocks are tables}
\label{sec:inputformat:tables}

A Yuzen case is a tree: named blocks, nested inside one another, each holding
keys and further blocks. That tree maps onto TOML in four steps.

\begin{description}
  \item[A block is a table.] \key{Grid} is written \code{[Grid]}, and every
    key after that header up to the next header belongs to it.
  \item[A nested block is a dotted table path.] \key{Evaporation} inside
    \key{Physics} is \code{[Physics.Evaporation]}; \key{SurfaceShape} inside
    \key{Initial} is \code{[Initial.SurfaceShape]}. The depth is not limited,
    and a container that carries no keys of its own ---
    \code{[Initial]}, \code{[Physics]} --- may be written as an empty table
    or left out entirely, since the dotted header of a child creates it.
  \item[A repeated block is an array of tables.] Write the header in double
    brackets, \code{[[Substrate.EllipseBump]]}, once per occurrence;
    \cref{sec:inputformat:repeat} says which blocks may repeat.
  \item[A vector-valued key is an array.] \code{Center = [0.5, 0.0]}. The two
    components are the $x$ and $y$ values, in that order, in 1-D runs as well.
\end{description}

Keys written before the first table header are the \textbf{root-level} keys:
\key{SolverName} and \key{JobName}, and the optional length-scale pair. TOML
requires them to come first, before any \code{[...]} line, because everything
after a header belongs to that table. A \key{JobName} placed after
\code{[Grid]} is a key of \key{Grid}, and the run stops for want of a
root-level \key{JobName}.

Block names, key names and string values are \textbf{case-sensitive}:
\key{Nx} is the grid size, and \code{nx} is an unknown key that stops the run
(\cref{sec:inputformat:unknown}). A string value that differs from an accepted
one only in case, such as \val{"ilu"} for \val{"ILU"}, is refused too, and the
message names the spelling that would have been accepted.

Key names are written in \code{PascalCase} (\key{NewtonMaxIterations},
\key{EndConcValue}), with American spelling (\key{Center}).

\section{How a written value becomes a number}
\label{sec:inputformat:types}

Conversion happens at read time and is decided by the type of the destination
variable in the solver, not by the TOML type you wrote. toml\texttt{++} hands
every scalar to \code{yuzen::config::Block} as text, and the text is converted
when the solver asks for the key. Writing \val{Nx = 300.0} is therefore legal,
and so is \val{Enable = 1}. The four conversions are strict: a value is converted whole or refused, and
the refusal names the key path and the value.

\begin{description}
  \item[bool] True for \val{true}, \val{yes}, \val{on}, \val{1} and
    \val{.true.}; false for \val{false}, \val{no}, \val{off}, \val{0} and
    \val{.false.}, compared case-insensitively. Anything else is refused.
  \item[int] Parsed as a double, which must be a whole number: \val{1e3} is
    1000, \val{1000.4} is refused. A whole number beyond the range of the
    key's integer type --- 32 bits, and 64 for \key{Numerics.MaxSteps} --- is
    clamped to its end.
  \item[double] Parsed with \code{strtod}; the whole text must be the number,
    and \val{nan} and \val{inf} are refused.
  \item[array] Converted element by element with the rules above. A list
    where one value is expected, a nested list or a TOML date or time is
    refused; so is nesting deeper than 64 levels.
\end{description}

The reader imposes no size on an array. Where a key is documented as an array
of two doubles, it is \file{Input.cpp} that checks the length afterwards and
aborts the run if it is not exactly~2. Beyond the conversion, values with
which the equations mean nothing --- a density, viscosity or surface tension
that is not positive, a time step, tolerance or output interval that is not
positive, a negative \key{EndTime}, an initial film with no liquid in it ---
are refused with \code{Input mismatch: key = value: why}; the reference
chapters give the ranges.

\begin{note}[What TOML catches for you]
  A number TOML cannot parse --- \val{1,0e-6} with a comma, \val{1e--6},
  \val{.5} without its leading zero --- is a \emph{syntax error}, and the run
  stops with the file name, the line and toml\texttt{++}'s description of the
  problem. The same goes for an unquoted string: \key{Dimension = xh} is not a
  value at all, where \key{Dimension = "xh"} is. A quoted value in a numeric
  key, such as \key{InitialDt = "fast"}, is valid TOML and is refused by the
  conversion above. Without strict conversions, \key{Enable =
  "ture"} would mean false, an unreadable double would become 0 and a date
  would be dropped as though the key had not been written;
  \file{tools/proptest.py} and the fuzzing of
  \cref{sec:verification:fuzzing} found these during development.
\end{note}

\section{Defaults, required keys and what is missing}
\label{sec:inputformat:defaults}

Whether a key is required is not declared anywhere in the file; it follows
from which overload of \code{getin} the solver calls for it.
\code{getin(name, v)}, with no default, throws \code{VariableNotFound}, which
\code{Input::Read} catches and turns into immediate termination of the
process. \code{getin(name, v, def)} stores \code{def} without a word. Blocks
work the same way: a block fetched with a bare \code{getBlockSize} or
\code{getBlock} is required, one fetched inside a
\code{try \dots{} catch(BlockNotFound)} is optional.

Thirteen keys are required outright across the whole file. They are \key{SolverName}
and \key{JobName} at root level, \key{Dimension}, \key{Nx}, \key{Xmin} and
\key{Xmax} in \key{Grid}, \key{Value} in \key{Initial.Concentration},
\key{Model} in \key{Physics.ConcField} (only when that block is present),
\key{InitialDt} in \key{Numerics}, \key{X} in every
\key{Output.ProbeList.Probe}, and \key{Start}, \key{End} and \key{PerDecade}
in \key{Output.FieldData.LogSchedule} (only when that block is present). Everything else has a built-in default. In the
reference chapters a required key is marked \reqd{} in the default column.

An optional block that is left out reads as an empty one: each of its keys
takes its default. The exceptions are in \key{Physics}: a missing sub-block
there is a term switched off, even where \key{Enable} defaults to
\val{true} inside a block that is written, and a missing
\key{Physics.ConcField} means \key{Model}~$=$~\val{"None"}
(\cref{ch:ref-physics}).

A required \emph{block} that appears twice is just as fatal as one that is
missing: the reader aborts with \code{Block [...] is required and must appear
exactly once}. In TOML
you cannot write a plain table header twice --- \code{[Grid]} twice is a
duplicate-key syntax error that toml\texttt{++} rejects --- so the way to trip
this is to write \code{[[Grid]]}, an array of tables, where a single table
belongs. The chapters that follow say for every block whether it is required
and whether it may repeat, and \cref{tab:file-map} summarizes it for the top
level.

\section{Entries the solver does not know stop the run}
\label{sec:inputformat:unknown}

Every key and every block the reader asks for is recorded, found or not. Once
the whole input has been read, anything in the file that nobody asked for is
something this version does not know --- a misspelled key, a key in the wrong
table, a table with a misspelled name --- and the run stops before anything is
computed, listing each one with its full path:

\begin{plain}
Error in the input: case.toml contains 3 entries that this version of Yuzen
does not read, so whatever they were meant to set would keep its default:
  key   Initial.PrecursorRatio   (this key is read in [Initial.SurfaceShape])
  key   Initial.SurfaceShape.EdgeBump.alpha   (did you mean Initial.SurfaceShape.EdgeBump.Alpha?)
  block [Physics.Evaporaton]   (did you mean [Physics.Evaporation]?)
Correct the spelling or remove the entry. The manual's input reference lists
every key each block accepts.
\end{plain}

A near miss of a name the block does read is offered as the likely intent; a
correctly spelled key in the wrong table is pointed at the table that reads it.
Like the other input errors, this one goes to standard error, before the result
directory exists (\cref{sec:output:running}).

To make these checks without running anything, give \code{--check-input}
before the file name:

\begin{shell}
yuzen --check-input case.toml
\end{shell}

The file is read and checked exactly as for a run --- unknown entries,
required keys, ranges, accepted values and the consistency checks between
keys --- and the program then prints \code{Input check passed} and exits
with status 0, or stops with the reason and status 1. Nothing is created.
A \key{Restart.FileName} that is not there, relative to the directory the
check is made in, is noted in a line starting \code{Note:} but does not fail
the check: the run that reads it may start elsewhere. A run started where it
is missing stops before making its result directory.

Without this, each of those entries would have been read by nobody and the
setting it was meant to change would have kept its default without a word:
the precursor at \val{0}, $\alpha$ at \val{1}, evaporation off. That is why
the check refuses rather than warns.

TOML's own checks are about syntax. A well-formed document with a key in the
wrong place is still well formed, which is why the check above is the solver's
and not toml\texttt{++}'s. What it cannot catch is a correctly spelled key, in
the right table, with the wrong value; for that, read back what the solver
understood --- the grid, the scales and the dimensionless groups it prints to
\textbf{standard output} (\cref{sec:output:running}).

\section{Repeated blocks}
\label{sec:inputformat:repeat}

Some blocks may appear more than once. These are the thirteen shape blocks,
seven under \key{Initial.SurfaceShape} and six under \key{Substrate}, together
with \key{Output.ProbeList.Probe}. The parser sizes the vector from
\code{getBlockSize} and fills it in file order, so four \key{EllipseBump}
blocks build four bumps. Zero copies are allowed as well --- a
\key{Substrate} table with no shape block in it is a flat substrate.

Write each occurrence as an array-of-tables header:

\begin{toml}
[[Substrate.RectBump]]
Order  = 1
Center = [0.0, 0.0]

[[Substrate.RectBump]]
Order  = 2
Center = [1.0e-3, 0.0]
\end{toml}

A block that may appear only once is a plain table, \code{[Substrate]}. The
shipped cases write a \emph{single} occurrence of a repeatable block as a plain
table too --- \code{[Substrate.EllipseBump]} --- and the solver reads it as a
list of one. Both spellings are accepted; use the double-bracket form as soon
as there is more than one, because a second plain \code{[Substrate.RectBump]}
header is a duplicate-key error rather than a second bump.

\section{Units}
\label{sec:inputformat:units}

\textbf{All dimensional input is SI}: meters, seconds, kilograms, \si{\pascal\second},
\si{\newton\per\metre}, \si{\metre\squared\per\second} and radians, except
where a table in the reference chapters says otherwise. Two exceptions are
worth memorizing now. The first is the contact angle,
\key{Physics.DisjoiningPressure.ContactAngle},
which is in \emph{degrees} --- and it is the only angle in the file that is;
the tilt of a \key{FlatPlane} shape, like everything else, is in radians. The
second is \key{Fluid.MolarMass}, which takes any consistent unit, because only
the ratio of its two values is ever used.

You never write a non-dimensional quantity, except for the handful of keys
that are genuinely dimensionless (exponents, ratios, switches and the
tolerances in \key{Numerics}). \code{NonDimensionalNumbers::Build} converts
your SI values into the solver's internal scales and prints the result to
standard output (\cref{sec:output:running}).

\section{A complete fragment}
\label{sec:inputformat:example}

\begin{toml}
# a 1-D film, 1 mm wide
SolverName = "Yuzen-1.0"
JobName    = "demo"

[Grid]
Dimension = "xh"
Nx        = 300
Xmin      = -0.5e-3
Xmax      = 0.5e-3

[Initial.Concentration]
Value = 0.5

[Initial.SurfaceShape]
TotalHeight    = 10e-6
PrecursorRatio = 0.02

[[Initial.SurfaceShape.EllipseBump]]
CreateType = "add"
Height     = 0.5
Center     = [0.0, 0.0]
Width      = [0.4e-3, 1.0]
SlopeWidth = 0.05

[Boundary.XminBC]
Type = "Symmetry"

[Boundary.XmaxBC]
Type = "Outlet"
\end{toml}

This fragment is not a complete input file: \key{Fluid}, \key{Physics},
\key{Output} and \key{Numerics} are required as well, and the run would abort
without them. Note that \code{[Initial]} and \code{[Boundary]} never appear as
headers of their own --- the dotted paths below them create both --- and that
the two root-level keys stand at the top, before any header.

\section{A map of the file}
\label{sec:inputformat:map}

\Cref{tab:file-map} lists everything that may appear at the top level of an
input file, and where in this part of the manual to find its keys. The
version covered throughout is Yuzen 1.0: 58 block paths, 54 of which carry
keys, and 229 key entries in total.

\begin{yzntable}{Every top-level block, and where it is documented\label{tab:file-map}}{@{}L{0.20\textwidth} L{0.13\textwidth} L{0.11\textwidth} L{0.44\textwidth}@{}}
\textsf{\bfseries Block} & \textsf{\bfseries Required} & \textsf{\bfseries Repeats} &
\textsf{\bfseries Contents, and where to read about it} \\ \midrule
Root level & partly & \dash & Keys before the first table header: \key{SolverName} and \key{JobName} are required, the length scale is optional. \Cref{sec:ref-domain:root}. \\
\key{Grid} & yes & no & Dimension, cell counts and domain bounds. \Cref{sec:ref-domain:grid}. \\
\key{Initial} & yes & no & Initial concentration and film shape, with repeatable shape sub-blocks. \Cref{sec:ref-domain:initial}. \\
\key{Substrate} & no & no & Substrate topography, with the same shape sub-blocks. \Cref{sec:ref-domain:substrate}. \\
\key{Fluid} & yes & no & Density, surface tension, viscosity, diffusivity, molar mass. \Cref{sec:ref-domain:fluid}. \\
\key{Physics} & yes & no & Which terms are active and their parameters. \Cref{ch:ref-physics}. \\
\key{Boundary} & yes & no & One condition per domain edge. \Cref{sec:ref-bcout:boundary}. \\
\key{Output} & yes & no & What is written and how often. \Cref{sec:ref-bcout:output}. \\
\key{Numerics} & yes & no & Time stepping, Newton and linear-solver settings. \Cref{sec:ref-numerics:numerics}. \\
\key{Restart} & no & no & Continue from a previous result. \Cref{sec:ref-numerics:restart}. \\
\end{yzntable}

No top-level block may be repeated; a second copy of a required one aborts the
run. The blocks that may repeat are all nested: the seven shape blocks of
\key{Initial.SurfaceShape}, the six of \key{Substrate}, and
\key{Output.ProbeList.Probe}.

\section{The \file{.inp} format}
\label{sec:inputformat:legacy}

Yuzen also reads a block-structured \file{.inp} format: any file name that
does not end in \file{.toml} goes to that parser. \textbf{The format is not
specified in this manual.} It is an alternative input format, and it is not
where new work belongs.

The shipped cases and every example in this manual are TOML.
